\documentclass{beamer}
\usepackage{xcolor}
\usepackage{amsmath}
\usepackage{amssymb}
\usepackage{graphicx}
\usepackage{booktabs}
% ── Footer ────────────────────────────────────────────────────────────────────
\setbeamertemplate{footline}{%
  \begin{beamercolorbox}[wd=\paperwidth, ht=4ex, dp=1ex, left, colsep=1ex]{frametitle}%
    \color{white}\textbf{Department of Biostatistics}%
    \hfill\color{white}\textbf{University of North Carolina at Chapel Hill}%
    \hfill\color{white}\textbf{Slide \insertpagenumber}%
  \end{beamercolorbox}%
}
\setbeamercolor{frametitle}{bg=blue!60!black, fg=white}
\setbeamertemplate{navigation symbols}{}
% ── Title ─────────────────────────────────────────────────────────────────────
\title{Group 8: Final Presentation\\
       {\normalsize Single-Word Speech Intelligibility for Children with Repaired Cleft Lip and Palate}}
\author{Esther Zhou \and Dan Huang \and Wanbing Wang \and Sophie Shan}
\date{
PI: Dr. David Zajac \\
Adams School of Dentistry \\
\today
}
\begin{document}
\frame{\titlepage}
% ══════════════════════════════════════════════════════════════════════════════
%  PART 1 – CONSULTATION PROCESS
% ══════════════════════════════════════════════════════════════════════════════
\begin{frame}
\begin{center}
  {\Large \textbf{\color{blue!60!black} Part I}}\\[1em]
  {\LARGE The Consultation Process}\\[0.5em]
  {\large How the project came together}
\end{center}
\end{frame}
% ── Slide 1: Initial Outreach ─────────────────────────────────────────────────
\begin{frame}
\frametitle{Initial Outreach: January 22, 2026}
\begin{itemize}
  \item Emailed \textbf{Dr.\ Preisser} for a collaborator introduction
  \item Introduced to \textbf{Dr.\ David Zajac}
  \item Dr.\ Zajac needed 1--2 weeks to organize the dataset
\end{itemize}
\end{frame}
% ── Slide 2: Meeting 1 ────────────────────────────────────────────────────────
\begin{frame}
\frametitle{Meeting \#1: February 20, 2026}
\begin{itemize}
  \item \textbf{Before:} Shared collaborator agreement and timeline; reviewed project overview; set agenda
  \item \textbf{During:} Project overview; analytic decisions; data codebook
  \item \textbf{After:} Received dataset; wrote SAP
\end{itemize}
\end{frame}
% ── Slide 3: Meeting 2 ────────────────────────────────────────────────────────
\begin{frame}
\frametitle{Meeting \#2: April 2, 2026}
\begin{itemize}
  \item \textbf{Before:} Drafted SAP; preliminary analysis; set agenda
  \item \textbf{During:} Reviewed SAP; explained LMM; previewed results
  \item \textbf{After:} Finalized analysis, tables, figures, and report
\end{itemize}
\end{frame}
% ── Slide 4: Challenges ───────────────────────────────────────────────────────
\begin{frame}
\frametitle{Challenges}
\begin{itemize}
  \item \textbf{Early scheduling delays:} Followed up proactively; sent materials in advance
  \item \textbf{Revised modeling plan:} Poor ICC in some listener groups undermined a mean SWI approach; switched to LMM on individual listener scores
\end{itemize}
\end{frame}
% ══════════════════════════════════════════════════════════════════════════════
%  PART 2 – THE ANALYSIS
% ══════════════════════════════════════════════════════════════════════════════
\begin{frame}
\begin{center}
  {\Large \textbf{\color{blue!60!black} Part II}}\\[1em]
  {\LARGE The Study}\\[0.5em]
  {\large Background, design, models, and results}
\end{center}
\end{frame}
% ── Slide: Study Background ───────────────────────────────────────────────────
\begin{frame}
\frametitle{Study Background}
\begin{columns}
  \column{0.6\textwidth}
  \textbf{Cleft Lip and Palate (CLP)}
  \begin{itemize}
    \item Common birth defect ($\sim$1 in 1{,}031 births)
    \item Structural differences affect speech airflow
    \item Reduced intelligibility may persist after surgical repair
  \end{itemize}
  \vspace{0.4em}
  \textbf{Speech Intelligibility:} how well a listener understands a speaker (a key clinical outcome)
  \column{0.38\textwidth}
  \includegraphics[width=\textwidth]{CLP.jpg}
\end{columns}
\end{frame}
% ── Slide: Research Aims ──────────────────────────────────────────────────────
\begin{frame}
\frametitle{Research Aims}
\begin{itemize}
  \item \textbf{Aim 1:} Assess differences in overall speech intelligibility between children with CLP and children in the CP, OM, and TD groups
  \item \textbf{Aim 2:} Evaluate whether patterns of phonetic errors differ between children with CLP and children in the CP, OM, and TD groups
\end{itemize}
\end{frame}
% ── Slide: Study Design ───────────────────────────────────────────────────────
\begin{frame}
\frametitle{Study Design}
\begin{itemize}
  \item Single-center, cross-sectional observational study; four groups: CLP, CP, TD, OM
  \item Children completed a 154-word speech test based on 11 phonetic contrasts
  \item Three adult listeners independently transcribed each child's recordings
  \item Outcome: \% words correctly identified $\to$ \textbf{3 correlated scores per child}
\end{itemize}
\end{frame}
% ── Slide: Data ───────────────────────────────────────────────────────────────
\begin{frame}
\frametitle{Data}
\textbf{Enrolled sample ($N=115$ children)}
\begin{itemize}
  \item CLP ($n=23$), CP ($n=24$), TD ($n=34$), OM ($n=34$)
  \item Mean age 7.6 years; 60\% male
  \item Mean intelligibility: 87.9\% (range: 44.2\%--98.1\%)
\end{itemize}
\vspace{0.5em}
\textbf{Analytic sample ($n=108$)}
\begin{itemize}
  \item 7 children with inadequate velopharyngeal (VP) status excluded (6 CLP, 1 CP)
  \item VP dysfunction is a physiological consequence of cleft palate; nearly fully confounded with CLP membership
  \item Full-sample analyses presented as sensitivity analyses
\end{itemize}
\end{frame}
% ── Slide: Preliminary Analysis ───────────────────────────────────────────────
\begin{frame}
\frametitle{Preliminary Analysis}
\textbf{1. Inter-rater reliability}
\begin{itemize}
  \item ICC assessed within each of 22 listener groups
  \item Listener Groups 2, 3, and 7 had poor agreement (ICC $< 0.5$)
  \item $\Rightarrow$ LMM on individual listener scores preferred over mean SWI
\end{itemize}
\vspace{0.4em}
\textbf{2. TD vs.\ OM comparability}
\begin{itemize}
  \item Bootstrapped TOST: observed difference = 3.08 points (90\% CI: 0.87, 5.68)
  \item Formal equivalence not confirmed, but difference clinically negligible ($\sim$3\% of scale)
  \item $\Rightarrow$ TD and OM merged into a single \textbf{Control} group
\end{itemize}
\vspace{0.4em}
\textbf{3. VP status}
\begin{itemize}
  \item CLP females with inadequate VP: mean SWI = 51.5 vs.\ 87.4 among adequate VP CLP females
  \item $\Rightarrow$ 7 children with inadequate VP excluded from primary and secondary analyses
\end{itemize}
\end{frame}
% ── Slide: Boxplot ────────────────────────────────────────────────────────────
\begin{frame}
\frametitle{Speech Intelligibility by Group and Sex}
\begin{center}
  \includegraphics[width=0.85\textwidth]{Figure1_4cohorts.pdf}
\end{center}
\end{frame}
% ── Slide: Scatter Plot ───────────────────────────────────────────────────────
\begin{frame}
\frametitle{Speech Intelligibility vs.\ Age by Cohort}
\begin{center}
  \includegraphics[width=0.75\textwidth]{Figure2_Scatter.pdf}
\end{center}
\end{frame}
% ── Aim 1 Intro ───────────────────────────────────────────────────────────────
\begin{frame}
\frametitle{Aim 1: Overall Speech Intelligibility Differences}
\begin{itemize}
  \item Assess differences in overall speech intelligibility between children with CLP and children in the CP and Control (OM + TD) groups
  \item Analytic sample: $n = 108$ (adequate VP status only)
\end{itemize}
\end{frame}
% ── Aim 1: Analytical Approach ────────────────────────────────────────────────
\begin{frame}
\frametitle{Aim 1: Analytical Approach}
\textbf{Method:} Linear mixed model (LMM)
\begin{itemize}
  \item Each child rated by 3 listeners $\to$ correlated observations
  \item Random intercept for child ID 
  \item Reference group: CLP
  \item Covariates: age, sex
  \item Pairwise interactions screened via likelihood ratio test (LRT)
  \item No interaction retained (all LRT $p \geq 0.40$); contrasts pooled across sex
  \item Multiplicity: Hochberg adjustment at $\alpha = 0.05$
\end{itemize}
\end{frame}
% ── Aim 1: LMM ────────────────────────────────────────────────────────────────
\begin{frame}
\frametitle{Aim 1: LMM}
\[
\text{SWI}_{ij} = \beta_0 + \beta_1\,\text{Cohort}_i + \beta_2\,\text{Age}_i + \beta_3\,\text{Sex}_i + u_{\text{ID}(i)} + \varepsilon_{ij}
\]
\vspace{0.2em}
\begin{itemize}
  \item No interaction terms retained (new\_Cohort:Age $p=0.677$; new\_Cohort:Sex $p=0.517$; Age:Sex $p=0.398$)
  \item $u_{\text{ID}(i)}$: child-level random intercept (SD = 5.97); listener group variance = 0
\end{itemize}
\end{frame}
% ── Aim 1: Adjusted Cohort Contrasts ─────────────────────────────────────────
\begin{frame}
\frametitle{Aim 1: Adjusted Cohort Contrasts}
\vspace{0.2em}
\renewcommand{\arraystretch}{1.35}
\begin{center}
\small
\begin{tabular}{@{}lrrr@{}}
\toprule
\textbf{Contrast} & \textbf{Estimate} & \textbf{95\% CI} & \textbf{Adjusted \(p\)} \\
\midrule
Control vs.\ CLP & $+$3.80 & ($-$0.07,\ 7.67) & 0.056 \\
CP vs.\ CLP      & $+$0.88 & ($-$3.71,\ 5.47) & 0.664 \\
\bottomrule
\end{tabular}
\end{center}
\normalsize
\vspace{0.5em}
\begin{itemize}
  \item Neither contrast reached statistical significance after Hochberg adjustment
  \item Control vs.\ CLP difference of 3.80 points approached significance ($p = 0.056$)
  \item Older age was positively associated with intelligibility ($p = 0.001$)
\end{itemize}
\end{frame}
% ── Aim 1: Fixed-Effect Estimates ────────────────────────────────────────────
\begin{frame}
\frametitle{Aim 1: Fixed-Effect Estimates}
\vspace{0.2em}
\scriptsize
\renewcommand{\arraystretch}{1.25}
\begin{center}
\begin{tabular}{@{}lccc@{}}
\toprule
\textbf{Term} & \textbf{Estimate} & \textbf{95\% CI} & \textbf{\(p\)-value} \\
\midrule
Intercept        & 70.18 & (60.32,\ 80.03) & $<0.001$ \\
Control vs.\ CLP & $+$3.80 & ($-$0.07,\ 7.67)\textsuperscript{$\dagger$} & 0.056   \\
CP vs.\ CLP      & $+$0.88 & ($-$3.71,\ 5.47)\textsuperscript{$\dagger$} & 0.664   \\
Age              & $+$2.07 & (0.84,\ 3.30)  & 0.001    \\
Sex (Female)     & $+$1.33 & ($-$1.09,\ 3.76) & 0.284  \\
\bottomrule
\end{tabular}
\end{center}
{\tiny $\dagger$ Hochberg-adjusted 95\% CI}
\normalsize
\vspace{0.2em}
\begin{itemize}
  \item Older age significantly associated with higher intelligibility ($p = 0.001$)
  \item Sex was not a significant predictor; no interaction terms retained
  \item Random intercept SD: child ID = 5.97; residual SD = 2.68
\end{itemize}
\end{frame}
% ── Aim 2 Intro ───────────────────────────────────────────────────────────────
\begin{frame}
\frametitle{Aim 2: Phonetic Error Patterns}
\begin{itemize}
  \item Evaluate whether patterns of phonetic errors differ between children with CLP and children in the CP and Control groups
  \item Same analytic sample ($n = 108$) and fixed effects as Aim 1
\end{itemize}
\end{frame}
% ── Aim 2: Analytical Approach ────────────────────────────────────────────────
\begin{frame}
\frametitle{Aim 2: Analytical Approach}
\textbf{Method:} 11 separate OLS regressions --- one per phonetic contrast
\begin{itemize}
  \item Same fixed effects as primary LMM (cohort, age, sex; no interactions retained)
  \item Reference group: CLP; contrasts pooled across sex
\end{itemize}
\vspace{0.4em}
\textbf{Multiple comparisons (Gatekeeping):}
\begin{itemize}
  \item \textbf{Family 1 (CLP vs.\ Control):} all 11 contrasts; Hochberg at $\alpha = 0.05$
  \item \textbf{Family 2 (CLP vs.\ CP):} only tested if Family 1 contrast is rejected
\end{itemize}
\end{frame}
% ── Aim 2: Regression Model ───────────────────────────────────────────────────
\begin{frame}
\frametitle{Aim 2: Regression Model}
\[
Y_{ik} =
\beta_0 +
\beta_1\,\text{Cohort}_i +
\beta_2\,\text{Age}_i +
\beta_3\,\text{Sex}_i +
\varepsilon_{ik}
\]
\begin{itemize}
  \item $Y_{ik}$: mean \% correct for contrast $k$, child $i$
  \item Fit separately for each of 11 phonetic contrasts
  \item No interaction terms; VP status not included (analytic sample: adequate VP only)
\end{itemize}
\end{frame}
% ── Aim 2: Full Results Table ─────────────────────────────────────────────────
\begin{frame}
\frametitle{Aim 2: All Phonetic Contrasts (Family 1)}
\vspace{0.1em}
\renewcommand{\arraystretch}{1.1}
\begin{center}
\tiny
\begin{tabular}{@{}lrrl@{}}
\toprule
\textbf{Phonetic Contrast} & \textbf{Estimate} & \textbf{95\% CI} & \textbf{Adj.\ \(p\)} \\
\midrule
Alveolar-palatal sibilant (see--she)     &  8.3 & (2.1,\ 14.5)    & 0.066 \\
Alveolar-velar stop (tap--cap)           &  2.7 & ($-$2.2,\ 7.7)  & 0.941 \\
Cluster-singleton (stop--top)            &  4.3 & ($-$1.1,\ 9.6)  & 0.923 \\
Initial consonant-null (pat--at)         &  1.7 & ($-$2.0,\ 5.3)  & 0.941 \\
Front-back vowel (fill--full)            &  3.0 & ($-$2.6,\ 8.5)  & 0.941 \\
Fricative-affricate (ship--chip)         &  5.0 & ($-$1.3,\ 11.4) & 0.923 \\
High-low vowel (beat--bit)               &  4.4 & ($-$0.2,\ 9.0)  & 0.591 \\
Liquid-glide (ray--way)                  &  7.0 & ($-$2.1,\ 16.2) & 0.923 \\
Stop-fricative/affricate (base--vase)    &  3.8 & ($-$2.0,\ 9.6)  & 0.941 \\
Stop-nasal (bop--mop)                    &  4.1 & ($-$0.0,\ 8.3)  & 0.511 \\
Voiced-voiceless stop (big--pig)         & $-$0.6 & ($-$5.7,\ 4.5) & 0.941 \\
\bottomrule
\end{tabular}
\end{center}
\normalsize
\vspace{0.2em}
{\small Estimates = Control $-$ CLP (mean \% correct), adjusted for age and sex. Family 2 not conducted for any contrast (all Family 1 adj.\ $p \geq 0.066$).}
\end{frame}
% ── Aim 2: Results ────────────────────────────────────────────────────────────
\begin{frame}
\frametitle{Aim 2: Results}
\textbf{No phonetic contrast was statistically significant after Hochberg adjustment}
\begin{itemize}
  \item All adjusted $p \geq 0.066$; Family 2 not conducted for any contrast
\end{itemize}
\vspace{0.4em}
\textbf{Largest estimated differences (Family 1, Control $-$ CLP):}
\renewcommand{\arraystretch}{1.2}
\begin{center}
\small
\begin{tabular}{@{}lrr@{}}
\toprule
\textbf{Phonetic Contrast} & \textbf{Estimate} & \textbf{95\% CI} \\
\midrule
Alveolar-palatal sibilant & 8.3 & (2.1,\ 14.5) \\
Liquid-glide              & 7.0 & ($-$2.1,\ 16.2) \\
Fricative-affricate       & 5.0 & ($-$1.3,\ 11.4) \\
\bottomrule
\end{tabular}
\end{center}
\normalsize
\vspace{0.2em}
{\small \textit{Estimates represent mean \% correct difference (Control $-$ CLP), adjusted for age and sex. None significant after Hochberg adjustment.}}
\end{frame}
% ── Slide: Conclusions ────────────────────────────────────────────────────────
\begin{frame}
\frametitle{Conclusions}
\begin{itemize}
  \item Neither the Control vs.\ CLP nor the CP vs.\ CLP contrast reached statistical significance; the Control vs.\ CLP difference of 3.80 points approached significance (adjusted $p = 0.056$)
  \item Older age was positively associated with intelligibility across all groups ($p = 0.001$)
  \item At the phonetic level, no contrast-specific difference survived Hochberg adjustment; Family 2 was not conducted
  \item Findings suggest children with repaired CLP may have lower speech intelligibility than controls, though larger studies are needed to draw definitive conclusions
\end{itemize}
\end{frame}
% ── Slide: Thank You ──────────────────────────────────────────────────────────
\begin{frame}
\frametitle{Thank You}
\begin{center}
  \vspace{0.5em}
  We thank \textbf{Dr.\ David Zajac} for his collaboration, patience, and generosity in sharing both his data and his clinical expertise.\\[0.5em]
  We also thank \textbf{Prof.\ John Preisser} for the introduction that made this project possible.\\[1em]
  \textbf{Questions?}\\[2em]
\end{center}
\end{frame}
\end{document}